% !TeX program = xelatex
% !TeX encoding = UTF-8
% 试题排版模板：使用 XeLaTeX 编译，替换页眉和正文占位文字即可。

\documentclass[a4paper,10pt]{article}
\usepackage[UTF8,fontset=fandol]{ctex} % 中文；Fandol 字体不依赖系统字体
\usepackage{amsmath,amssymb}          % 数学公式与符号
\usepackage{graphicx}                % 插图
\usepackage{fancyhdr}                % 页眉、页脚
\usepackage{enumitem}                % 题号与题间距
\usepackage{tabularx}                % 选择题选项列宽

% 页面尺寸：沿用原稿的紧凑布局，按需调整各边距。
\usepackage[
    a4paper,
    top=0.5cm,
    bottom=0.8cm,
    left=1.7cm,
    right=1.7cm,
    headheight=20pt,
    headsep=0.5cm,     % 页眉与正文的距离
    footskip=0.5cm,    % 页码与正文的距离
    includehead,
    includefoot
]{geometry}
\setlength{\parindent}{0pt}

% 页眉信息：修改以下三项；文字尽量简短，避免重叠。
\pagestyle{fancy}
\fancyhf{}
\fancyhead[L]{\small 20XX 年第 X 学期}
\fancyhead[C]{\small 学校／学院}
\fancyhead[R]{\small 课程名称（考试类型）}
\renewcommand{\headrulewidth}{1pt}
\fancyfoot[C]{\thepage}

% 大题标题：统一字号及标题前后的间距。
\newcommand{\examsection}[1]{%
    \par\addvspace{0.3cm}%
    {\Large #1\par}%
    \nopagebreak[4]\vspace{0.4cm}%
}

% 每个 enumerate 环境内用 \item 添加题目，题号从 1 开始。
\setlist[enumerate,1]{
    label=\textbf{\arabic*.},
    leftmargin=1.8em,
    labelwidth=1.4em,
    labelsep=0.4em,
    align=left,
    topsep=0pt,
    partopsep=0pt,
    parsep=0pt,
    itemsep=0.52cm     % 相邻题目之间的距离
}

\begin{document}

\examsection{一、单项选择题（每题 X 分，共 X 题）}
\begin{enumerate}
    \item 【选择题题干】（\qquad）
    % 下表为一行四个选项；选项较长时可改为逐行排列。
    \par\vspace{0.2cm}
    \begin{tabularx}{\linewidth}{@{}*{4}{>{\raggedright\arraybackslash}X}@{}}
        A. 【选项】 & B. 【选项】 & C. 【选项】 & D. 【选项】
    \end{tabularx}
\end{enumerate}

\examsection{二、填空题（每题 X 分，共 X 题）}
\begin{enumerate}[itemsep=0.2cm]
    \item 【填空题题干】\underline{\hspace{3cm}}。
    % 修改上面的 3cm 可调整填空横线长度。
\end{enumerate}

\examsection{三、论述题（每题 X 分，共 X 题）}
\begin{enumerate}
    \item 【论述题题干】
    \par\vspace*{3cm} % 答题留白，按需修改高度
\end{enumerate}

\examsection{四、作图题（每题 X 分，共 X 题）}
\begin{enumerate}
    \item 【作图题题干】
    % 插图时取消下面三行注释，并替换图片路径；路径相对于本文件。
    % \begin{center}
    %     \includegraphics[width=0.65\linewidth]{题图.png}
    % \end{center}
    \par\vspace*{4cm}
\end{enumerate}

\examsection{五、计算题（每题 X 分，共 X 题）}
\begin{enumerate}
    \item 【计算题题干】
    \begin{enumerate}[label=(\arabic*),leftmargin=2em,itemsep=0.15cm]
        \item 【第一个小问】（X 分）
        \item 【第二个小问】（X 分）
    \end{enumerate}
    \par\vspace*{5cm}
\end{enumerate}

% 需要另起一页时，在相应题目或大题前插入 \newpage。
\end{document}
